% Draft abstract and introduction for abstract.tex (PLanQC 2027).
% To use: replace the abstract environment of abstract.tex with the one below, and
% replace Section 1 ("The problem", from \section{The problem} up to
% \section{Separating ...}) with the section below. It uses the macros and labels
% already defined in preamble.tex, statements.tex and the later sections.

\begin{abstract}
  Quantum programming languages let a program discard a temporary register under sufficient conditions: a
  syntactic restriction on how the temporary was produced, a required final state, or the programmer's promise.
  We give the exact semantic condition behind these rules. A drop is \emph{safe} when the program prefix that
  ends at it, with every register exposed, is observationally equivalent to a program that prepares the
  temporary in a fixed state, independently of the input, and leaves it idle: up to observational equivalence,
  the drop is a weakening. We analyze this definition into a decidable criterion. Without measurement,
  observational equivalence of programs that return every register is equality of denotations up to a global
  phase, and the programs that prepare the temporary independently denote exactly the maps $g\otimes W$, with $g$
  a unit vector on the temporary and $W$ an isometry on the other registers. A drop is therefore safe exactly
  when the prefix's denotation $K$ factors as $g\otimes W$, that is, when the slices $K_i=(\bra i\otimes I)K$
  along the temporary's computational basis, the Kraus operators of the channel the drop induces on the other
  registers, are pairwise proportional. With measurement, definition and criterion apply to the unnormalized map
  of each measurement record, so the verdict can depend on the outcome. The criterion extends to borrowed
  registers, and deciding it is coNP-complete for programs over $X$, CNOT and Toffoli, so no sound
  polynomial-time rule is complete for them unless $\mathrm P=\mathrm{NP}$. We implement it in a checker and a compiler to
  OpenQASM~3.
\end{abstract}

\section{Introduction}\label{sec:intro}\label{sec:problem}

Quantum programming languages let a program discard a temporary register when its scope ends, and each guards
the discard with its own rule. Silq and Unqomp uncompute temporaries whose producers are
\code{qfree}~\cite{Bichsel2020Silq,Paradis2021Unqomp}; Q\#, Quipper, Qutes and ReQWIRE require the temporary to
end in a designated state, on the programmer's word or by
analysis~\cite{Svore2018Qsharp,MicrosoftQsharpMemory,Green2013Quipper,Faro2026Qutes,Rand2019ReQWIRE}, and
fractional types require the same of reversible programs~\cite{Chen2020Fractional}. Each rule is a sufficient
condition for a single semantic property: the program behaves as if the temporary had been prepared on its own
and left idle. We define this property by observational equivalence and analyze the definition into an exact
criterion that a compiler can check. Exactness has a price: deciding the criterion is coNP-complete for programs
over $X$, CNOT and Toffoli, so every sound polynomial-time rule rejects some safe drops unless
$\mathrm P=\mathrm{NP}$.

We write $\Omega\vdash M:A$ for a program with input registers $\Omega$ and output type $A$. Without measurement
it denotes an isometry $\den M:\den\Omega\to\Hilb_{\mathsf{hid}}\otimes\den A$, where $\Hilb_{\mathsf{hid}}$ holds
the registers $M$ has dropped: \drop{} ends a register's scope and leaves the state unchanged. A closed context
$C$ supplies the inputs of a program and may compute further on its outputs; the observable result of a closed
program is the probability of each measurement record and of each standard-basis outcome of its output. Two
programs of the same type are \emph{observationally equivalent}, $M\cong M'$, when $C[M]$ and $C[M']$ have the same
observable result for every closed context $C$. Now fix a drop site $s$ in $M$ whose block $v$ is a tuple of
temporaries allocated in the state $\psi$, and let $P_1=\pfx{M}{s}$ be the prefix of $M$ before $s$ with every
\drop{} deleted, so that $P_1$ returns every register; a context observes a register's correlations only while
the register is in scope. The drop is \emph{safe} when $P_1\cong P_2$ for some \emph{witness}
\[
  P_2\;=\;\mathtt{new}\ v\ \mathtt{@}\ \psi\mathtt{;}\ \ R\ v\mathtt{;}\ \ Q,
\]
with $R$ a unitary on the block and $Q$ a program that acts only on the other registers. In $P_2$ the block
is prepared on its own and left idle, so dropping it is a weakening: safety says that the drop is a weakening up
to observational equivalence. Without measurement, programs that return every register are observationally
equivalent exactly when their denotations agree up to a global phase, which $Q$ can absorb. Safety thus asks for
\[
  \den{P_1}\;=\;\den{P_2}.
\]

\Cref{ex:three} gives three well-typed programs. Each allocates $t$ in $\ket0$, computes, and drops $t$ while
returning $v$; purple marks the lines that (b) and (c) add to (a).

\begingroup
\setlength{\topsep}{2pt}
\exRunning
\endgroup

Program~(a) is safe: its prefix already has the shape of $P_2$, with $R=H$. In program~(b), the \code{CX}
entangles $t$ with $v$ at input $\ket0_v$, whereas every $P_2$ returns $t$ in the product state $R\psi$, so its
drop is unsafe. The temporary is dead after the drop in both programs; scope information alone cannot decide
safety. Program~(c), treated in \cref{sec:results}, is safe on one measurement outcome and unsafe on the other.
The definition is semantic: the prefix of Deutsch--Jozsa applies the oracle to its temporary, yet it is
observationally equivalent to a $P_2$ that prepares the temporary in $\ket-$ and applies the phase oracle to the
query register alone.

The definition quantifies over all contexts and all witnesses $P_2$; analyzing it reduces safety to a property
of one linear map. Let $K=\den{P_1}:\den\Omega\to\Hv\otimes\Hrest$, where $\Hrest$ holds every register outside
the block, including those dropped earlier. Unfolding the witness gives $\den{P_2}x=R\psi\otimes\den Qx$, so
safety asks for $K=g\otimes W$ with $g=R\psi$ and $W=\den Q$. Conversely, since our gates include every unitary,
each unit vector $g$ is $R\psi$ for some $R$, and each isometry $W$ is $\den Q$ for some $Q$ that extends $W$ to
a unitary. Expanding $K$ along the computational basis of the block, with slices $K_i=(\bra i\otimes I)K$, then
gives
\[
  \text{the drop is safe}
  \iff K=g\otimes W
  \iff K_i=g_i\,W\ \text{for every } i
  \iff \operatorname{rank}\Gamma\le1,
\]
where $\Gamma_{ij}=\operatorname{tr}(K_i^{\dagger}K_j)$ is the Gram matrix of the slices. The first equivalence
is \cref{thm:measurementfree}; the others are linear algebra. The slices are the Kraus operators of the channel
$\rho\mapsto\sum_iK_i\rho K_i^{\dagger}$ that the drop induces on the other registers, so a drop is safe exactly
when this channel has Kraus rank one. The unit vector $g$, the \emph{certificate}, is the state in which the
block ends at every input. Program~(a) has $K_0=K_1=I/\sqrt2$ and $g=\ket+$; program~(b) has $K_0=I/\sqrt2$ and
$K_1=X/\sqrt2$, which are linearly independent. The criterion accepts (a), whose producer $H$ lies outside
\code{qfree}, and the temporary of Deutsch--Jozsa, which ends in $\ket-$ for every oracle, whereas designated-state
rules require $\ket0$.

With measurement, a program denotes one unnormalized map per measurement record, and observational equivalence
compares the probability of each record as well. Safety is defined record by record, against witnesses that may
realize the remaining computation with one extra qubit and a measurement of which they keep one outcome. The
same analysis shows that the drop is safe on a record $o$ exactly when the prefix's map on $o$ factors as
$\sigma_s^o=g\otimes W_o$ with $W_o$ a contraction, where $\lVert W_ox\rVert^2$ is the probability of $o$ at a
unit input $x$ (\cref{thm:perrecord}). Program~(c) is safe on exactly one of its
two records. Coarser semantics cannot decide safety: two programs can share their output channel and record
probabilities while receiving opposite verdicts on every record (\cref{sec:status}).

For borrowed registers, which must return to the caller with their entanglement intact, the certificate
generalizes to an isometry $B$ on the borrowed inputs, $K=B\otimes W$. A safe block is in a known state, so a
compiler can reclaim it, and reusing reclaimed temporaries preserves every record, its probability and the
reduced output state. We prove type safety and these results for a typed first-order language with
finite-dimensional registers, measurement and classical conditionals, and implement the criterion in Haskell,
with exact arithmetic for the verdicts and an OpenQASM~3 back end that emits reclaiming gates
(\cref{sec:status}).
